\documentclass[12pt]{article}

\usepackage{amsmath, amssymb, amsthm}
\usepackage{geometry}
\usepackage{graphicx}
\usepackage{booktabs}
\usepackage{tikz}
\usepackage{hyperref}
\usepackage{float}
\usepackage{xcolor}
\usepackage{tikz}
\usepackage{needspace}
\usepackage{placeins}
\usetikzlibrary{arrows.meta}
\numberwithin{equation}{section}
\usepackage{hyperref}
\usepackage{url}
\usepackage{caption}
\usepackage{listings}
\captionsetup[listing]{
	skip=4pt
}
\usepackage{fancyhdr}

\pagestyle{fancy}
\fancyhf{}

\fancyhead[L]{\nouppercase{\rightmark}}
\fancyhead[R]{\thepage}

\renewcommand{\subsectionmark}[1]{%
	\markright{\thesubsection\quad #1}
}

\title{
	\textbf{European Option Convergence Study}\\[0.0em]
}

\author{Neel Dev Sen}
\date{August 09, 2026}
\author{Neel Dev Sen}
\usepackage[newfloat]{minted}
\usemintedstyle{manni}

\begin{document}
    \maketitle
	\newpage
	\tableofcontents
	\newpage
	
	\newpage
	\setcounter{section}{0}
	
	
\section*{Introduction}
This convergence study is based on the 3 European Option pricers I have made in my other project called \textbf{Pricing Engines} (which you can find on my \href{https://github.com/neeldevsen/pricing-engines/tree/main}{\textbf{GitHub}}). The numerical methods that will be compared is the Binomial Tree, Trinomial Tree and the Finite Differences (using the Crank-Nicolson scheme)

\section{Theoretical  and Numerical Framework}
This section will mainly just be stating the Time and Space complexities, the complete derivations of these are in this \href{https://github.com/neeldevsen/pricing-engines/blob/main/PricingEngines.pdf}{\textbf{document}}. 
\subsection{Time Complexities}
The time complexity of a Binomial Tree with $N$ steps is:
\begin{equation}
	\boxed{\mathcal{O}(N^2)}
\end{equation}
The time complexity of a Trinomial Tree with $N$ steps is also:
\begin{equation}
	\boxed{\mathcal{O}(N^2)}
\end{equation}
Finally the time complexity of the Finite Differences with a mesh of $M$ space nodes and $N$ time nodes is:
\begin{equation}
	\boxed{\mathcal{O}(M\times N)}
\end{equation}
Since I used the Thomas algorithm it makes it faster instead of $\mathcal{O}(M^3\times N)$
In this experiment to keep this all similar so it is easier to compare, I define $M=N$ to make a square mesh such that the time complexity of Finite Differences is also:
\begin{equation}
	\boxed{\mathcal{O}(N^2)}
\end{equation}
All of these time complexities will be useful when I show the runtime-convergence graphs of these numerical methods
\subsection{Space Complexities}
The space complexity of a Binomial Tree with $N$ steps is:
\begin{equation}
	\boxed{\mathcal{O}(N)}
\end{equation}
The space complexity of a Trinomial Tree with $N$ steps is also:
\begin{equation}
	\boxed{\mathcal{O}(N^2)}
\end{equation}
Finally the time complexity of the Finite Differences with a mesh of $M$ space nodes and $N$ time nodes is:
\begin{equation}
	\boxed{\mathcal{O}(M)}
\end{equation}
Over here I don't store the entire matrix, instead I store the 3 diagonals as 3 separate vectors. Since again I am using $M=N$ this space complexity is equal to:
\begin{equation}
	\boxed{\mathcal{O}(N)}
\end{equation}
This will be useful for the section on memory. 
\subsection{Error Metrics}
The \textbf{absolute error} is the main error metric that will be used to show convergence and it goes as shown.
\begin{equation}
	\boxed{\operatorname{Absolute  \hspace{4pt} Error} =  |V_N - V_{CF} |}
\end{equation}
where $V_N$ is the value of the option given by the numerical method, and $V_{CF}$ is the value of the option given by the closed form solution. 
\newline 
For the Arbitrage Error, I will just look at how far the options are from meeting the \textbf{Put-Call Parity} (the no-arbitrage condition)
The Put Call Parity is as shown
\begin{equation}
	C_t - P_t = S - Ke^{-r(T-t)}
\end{equation}
The arbitrage error will be defined as:
\begin{equation}
	\boxed{\operatorname{Arbitrage \hspace{4pt} Error} = |C_t - P_t - S + Ke^{-r(T-t)}|}]
\end{equation}
\newpage
\subsection{Data}
The table below will be the data that is used in the convergence study. 
\begin{table}[H]
	\centering
	\begin{tabular}{lccc}
		\hline
		\textbf{Parameter} & \textbf{Black--Scholes} & \textbf{Black--76} & \textbf{Merton} \\
		\hline
		Spot Price $S_0$       & 120.0 & --    & 120.0 \\
		Futures Price $F_0$    & --    & 115.0 & --    \\
		Strike $K$             & 100.0 & 100.0 & 100.0 \\
		Risk-Free Rate $r$     & 0.04  & 0.04  & 0.04  \\
		Dividend Yield $q$     & --    & --    & 0.02  \\
		Volatility $\sigma$    & 0.25  & 0.25  & 0.25  \\
		Maturity $T$           & 2.0   & 2.0   & 2.0   \\
		Domain Multiplier      & 4.0   & 4.0   & 4.0   \\
		\hline
	\end{tabular}
	\caption{Parameters used for the numerical convergence study.}
	\label{tab:parameters}
\end{table}

\section{Results: Black-Scholes}
\subsection{Black-Scholes Call}
\subsection{Black-Scholes Put}
\subsection{Arbitrage Error}
\subsection{Runtime}

\section{Results: Black-76 }
\subsection{Black-76 Call}
\subsection{Black-76 Put}
\subsection{Arbitrage Error}
\subsection{Runtime}

\section{Results: Merton}
\subsection{Merton Call}
\subsection{Merton Put}
\subsection{Arbitrage Error}
\subsection{Runtime}

\section{Conclusion}


\end{document}